\documentclass[aspectratio=169,12pt]{beamer}
% For a study PDF with quiz answers visible, add handout to the class options.
\usepackage{nulifeinsurance}
\NUSetup{Data Science for Life Insurance}{Week 2}{nu-wordmark.png}
\title{Business Understanding and Requirements}
\subtitle{Life Insurance Policy Retention\\Understanding and Anticipating Lapse}
\author{Instructor Name}
\institute{National University}
\date{}
\begin{document}
\begin{frame}[plain]\titlepage\end{frame}
\begin{frame}{The business situation}
\begin{insuranceexample}
Illustrative capacity: 100 policy reviews each month. Confirm customer-contact capacity separately.
Which active policies should receive support, and what evidence is needed?
\end{insuranceexample}
\begin{conceptblock}
Begin with the decision, intended user, available action, and constraints.
\end{conceptblock}
\end{frame}
\begin{frame}{Where are we in the life cycle?}
\NULifecycle{\textbf{Business understanding} $\rightarrow$ Data understanding
$\rightarrow$ Preparation $\rightarrow$ Analysis and modeling
$\rightarrow$ Evaluation $\rightarrow$ Delivery and monitoring}
\begin{commonmistake}
``Build a classifier'' describes a technical approach, not the business requirement.
\end{commonmistake}
\end{frame}
\NUQuickQuiz{Which statement is a business requirement?}
{\begin{enumerate}
\item Use a decision tree.
\item Prioritize 100 active policies for monthly retention review.
\end{enumerate}}
{Statement 2 specifies a decision and an operational constraint.
The model is an implementation choice.}
\begin{frame}{Your project mission}
\begin{projecttask}
Interview the retention manager and write a one-page project charter.
\end{projecttask}
\begin{itemize}
\item Define lapse and distinguish other policy exits.
\item Agree on the prediction horizon and outreach capacity.
\item Record success criteria, limitations, and unanswered questions.
\end{itemize}
\end{frame}
\end{document}
